\usepackage[a4paper,left=10mm,right=10mm,top=9mm,bottom=12mm,footskip=6mm]{geometry}
\usepackage[T1]{fontenc}
\usepackage[utf8]{inputenc}
\usepackage[ngerman]{babel}
\usepackage{lmodern,microtype,amsmath,array,tabularx,xcolor,listings,tcolorbox,fancyhdr,hyperref}
\renewcommand{\familydefault}{\sfdefault}
\definecolor{ink}{HTML}{182C3A}
\definecolor{accent}{HTML}{294F66}
\definecolor{string}{HTML}{345B39}
\definecolor{muted}{HTML}{50545A}
\definecolor{paper}{HTML}{F3F4F5}
\definecolor{rule}{HTML}{A7ADB1}
\color{ink}
\hypersetup{colorlinks=true,urlcolor=accent,pdftitle={Python-Kurzreferenz | Digitalisierung in Unternehmen},pdfauthor={Dr. Christopher Gerling}}
\setlength{\parindent}{0pt}
\setlength{\parskip}{.7pt}
\setlength{\tabcolsep}{3pt}
\renewcommand{\arraystretch}{1.13}
\pagestyle{fancy}\fancyhf{}
\renewcommand{\headrulewidth}{0pt}
\fancyfoot[L]{\fontsize{6.3}{7.5}\selectfont Dr. Christopher Gerling\enspace·\enspace HWR Berlin\enspace·\enspace Digitalisierung in Unternehmen}
\fancyfoot[R]{\fontsize{6.3}{7.5}\selectfont Stand 17.11.2026\enspace |\enspace\thepage\,/\,2}
\lstset{language=Python,basicstyle=\ttfamily\fontsize{7.7}{8.8}\selectfont\color{ink},keywordstyle=\bfseries\color{accent},stringstyle=\color{string},commentstyle=\itshape\color{muted},showstringspaces=false,keepspaces=true,columns=fixed,basewidth=4.04pt,upquote=true,aboveskip=2pt,belowskip=2pt,tabsize=4,breaklines=false,morekeywords={True,False,None},literate={ä}{{\"a}}1 {ö}{{\"o}}1 {ü}{{\"u}}1 {Ä}{{\"A}}1 {Ö}{{\"O}}1 {Ü}{{\"U}}1 {ß}{{\ss}}1}
\newcommand{\py}[1]{\mbox{\ttfamily\fontsize{7.7}{8.8}\selectfont\detokenize{#1}}}
\newtcolorbox{karte}[2][]{before upper=\raggedright,colback=white,colframe=rule,colbacktitle=paper,coltitle=ink,fonttitle=\bfseries\fontsize{9.4}{11}\selectfont,title={#2},boxrule=.45pt,arc=1mm,left=2mm,right=2mm,top=.4mm,bottom=.5mm,toptitle=.8mm,bottomtitle=.8mm,before skip=0pt,after skip=.7mm,#1}
\newcommand{\merk}[1]{\par{\bfseries Merke:} #1\par}
\newenvironment{tabelle}[1]{\tabularx{\linewidth}{@{}#1@{}}}{\endtabularx}
% Beide Kopfzeilen nutzen gemeinsame Grundlinien links und rechts.
\newcommand{\kopf}{%
\noindent\hbox to \linewidth{{\fontsize{19}{21}\selectfont\bfseries Python 3\enspace\mdseries Kurzreferenz}\hfil{\fontsize{8.1}{9.4}\selectfont\bfseries Dr. Christopher Gerling}}\par
\vspace{1mm}
\noindent\hbox to \linewidth{{\fontsize{9.5}{12}\selectfont Digitalisierung in Unternehmen · BWL Bank}\hfil{\fontsize{8.1}{9.4}\selectfont HWR Berlin}}\par
\vspace{2mm}{\color{accent}\hrule height 1pt}\vspace{3mm}}
